\documentclass[a4paper,11pt]{article}
\usepackage{exam-style-341}
\examinehead{341 农业综合知识三 · 数据库技术与应用}{模拟试卷（二）参考答案}
\begin{document}
\answercover{341 农业综合知识三 · 数据库技术与应用}{模拟试卷（二）参考答案}

本答案按试卷题号顺序给出，每部分均标注分值。选择题、SQL 题与规范化题附有解析；规范化题给出闭包与候选码推导、范式判定、分解过程，并用表格法（Chase 算法）验证无损连接性。

% ============================================================
\answerpart{一、单项选择题（10 题，每题 2 分，共 20 分）}
% ============================================================

\q \textbf{C}。
\begin{analysis}
内模式（也称存储模式）是数据库在物理存储方面的描述，是数据物理结构和存储方式的描述，一个数据库只有一个内模式；外模式是用户能够看见和使用的局部数据的逻辑结构；模式是全体数据的全局逻辑结构；视图是从基本表导出的虚表，属于外模式范畴。故选 C。
\end{analysis}

\q \textbf{B}。
\begin{analysis}
模式/内模式映像定义数据库全局逻辑结构与物理存储结构之间的对应关系。当数据库的物理存储结构改变时，DBA 修改该映像可使模式保持不变，从而应用程序也不必改变，这就是数据的物理独立性。外模式/模式映像保证的是数据的逻辑独立性。故选 B。
\end{analysis}

\q \textbf{B}。
\begin{analysis}
数据模型的三要素中，数据结构描述系统的静态特性（数据库的组成对象及其联系）；数据操作描述系统的动态特性，即对数据库中各种对象允许执行的操作集合以及有关的操作规则；完整性约束是一组完整性规则，用以保证数据的正确、有效和相容。故选 B。
\end{analysis}

\q \textbf{C}。
\begin{analysis}
若关系 $R$ 中的属性（组）$X$ 不是 $R$ 的主码，但 $X$ 是另一个关系 $S$ 的主码，则称 $X$ 是 $R$ 的外码（外键）。候选码是 $R$ 中能唯一标识元组的属性组；主属性是包含在任一候选码中的属性；超码是能唯一标识元组但可能含多余属性的属性组。故选 C。
\end{analysis}

\q \textbf{B}。
\begin{analysis}
参照完整性规则：若属性（组）$F$ 是基本关系 $R$ 的外码，它与基本关系 $S$ 的主码 $K$ 相对应，则 $R$ 中每个元组在 $F$ 上的值必须或者取空值，或者等于 $S$ 中某个元组的主码值。选项 A 是实体完整性的内容；选项 C 是第一范式（属性不可再分）的要求；选项 D 是关系的性质之一（关系中不允许出现完全相同的元组）。故选 B。
\end{analysis}

\q \textbf{C}。
\begin{analysis}
各级范式消除的函数依赖问题如下：1NF 消除非原子属性（属性不可再分）；2NF 在 1NF 基础上消除非主属性对码的部分函数依赖；3NF 在 2NF 基础上消除非主属性对码的传递函数依赖；BCNF 在 3NF 基础上进一步消除主属性对码的部分函数依赖和传递函数依赖。多值依赖是 4NF 所消除的。因此与 3NF 相比 BCNF 进一步消除了“主属性对码的部分函数依赖和传递函数依赖”，选 C。此外，若 $R$ 属于 BCNF，则 $R$ 一定属于 3NF，反之不成立。
\end{analysis}

\q \textbf{B}。
\begin{analysis}
自然连接要求两个关系中进行比较的分量必须是同名属性组，并在结果中去掉重复的属性列。$R(A,B,C)$ 与 $S(B,C,D)$ 的公共属性组为 $(B,C)$，连接条件为 $R.B=S.B$ 且 $R.C=S.C$，结果关系模式为 $(A,B,C,D)$，共 4 个属性。
\end{analysis}

\q \textbf{A}。
\begin{analysis}
原子性（Atomicity）指事务是数据库的逻辑工作单位，事务中包括的所有操作要么都做，要么都不做。选项 B 描述持久性；选项 C 描述一致性；选项 D 描述隔离性。故选 A。
\end{analysis}

\q \textbf{C}。
\begin{analysis}
视图是从一个或几个基本表（或视图）导出的表，它是一个虚表，数据库中只存放视图的定义而不存放视图对应的数据，视图中的数据仍存放在导出它的基本表中，所以选项 C 错误。视图可以嵌套定义；视图能对机密数据提供安全保护（用户只能通过视图访问权限内的数据）；定义视图后对视图的操作最终转换为对基本表的操作，因此视图为数据库的重构提供了一定程度的逻辑独立性。故选 C。
\end{analysis}

\q \textbf{A}。
\begin{analysis}
两段锁（2PL）协议规定：所有事务必须分两个阶段对数据项加锁和解锁——扩展阶段只能加锁不能解锁，收缩阶段只能解锁不能加锁。可以证明，若并发事务都遵守两段锁协议，则这些事务的任何并发调度都是可串行化的，因此遵守两段锁协议是可串行化调度的\textbf{充分条件}（但不是必要条件，可串行化调度不一定遵守两段锁协议），A 正确。遵守两段锁协议的事务仍可能发生死锁（如两个事务互相等待对方持有的锁），B 错误。封锁粒度越大，被封锁的数据范围越大，并发度越低，C 错误。共享锁（S 锁）只允许其他事务读取，不允许修改，D 错误。
\end{analysis}

% ============================================================
\answerpart{二、填空题（5 题，每题 2 分，共 10 分）}
% ============================================================

\q 第一空：\textbf{同名（相同）}；第二空：\textbf{重复的属性列}。
\begin{analysis}
自然连接是一种特殊的等值连接：它要求两个关系中进行比较的分量必须是同名属性组（公共属性），并且在结果中把重复的属性列去掉，只保留一份。
\end{analysis}

\q 第一空：\textbf{主属性}；第二空：\textbf{非主属性}。
\begin{analysis}
候选码是能唯一标识元组且不含多余属性的属性组。包含在任何一个候选码中的属性称为主属性（码属性）；不包含在任何候选码中的属性称为非主属性（非码属性）。
\end{analysis}

\q \textbf{用户定义}。
\begin{analysis}
关系数据库的三类完整性约束是实体完整性、参照完整性和用户定义完整性。其中实体完整性和参照完整性是关系模型必须满足的完整性约束，由 DBMS 自动支持；用户定义完整性则是针对某一具体应用环境而由用户自行定义的约束条件（如“成绩必须在 0～100 之间”）。
\end{analysis}

\q 第一空：\textbf{$ABCD$}；第二空：\textbf{$A$}。
\begin{analysis}
（1）$(AC)^{+}$：初始为 $\{A,C\}$；由 $A\to B$ 得 $\{A,B,C\}$；由 $B\to C$ 未产生新属性；由 $C\to D$ 得 $\{A,B,C,D\}$。故 $(AC)^{+}=ABCD$。
（2）$A^{+}$：初始为 $\{A\}$；由 $A\to B$ 得 $\{A,B\}$；由 $B\to C$ 得 $\{A,B,C\}$；由 $C\to D$ 得 $\{A,B,C,D\}$，即 $A^{+}$ 包含了 $R$ 的全部属性，故 $A$ 是候选码。又 $A$ 是单属性，不能再去掉任何属性，故 $R$ 的候选码就是 $A$（唯一）。
\end{analysis}

\q \textbf{非过程性}。
\begin{analysis}
过程性语言要求用户既要指出“做什么”，又要指出“怎么做”（如按什么顺序、用什么方法访问数据）；非过程性语言只需指出“做什么”，不必指出“怎么做”，具体存取路径与执行过程由 DBMS 的优化器和执行器完成。SQL 属于非过程性（描述性）语言，这正是 SQL 简单易学、使用方便的重要原因。
\end{analysis}

% ============================================================
\answerpart{三、简答题（2 题，每题 5 分，共 10 分）}
% ============================================================

\q E-R 模型及其作用、联系的类型：
\begin{analysis}
\textbf{（1）什么是 E-R 模型（2 分）}

E-R 模型（Entity-Relationship Model，实体-联系模型）是用 E-R 图来描述现实世界概念结构的数据模型，属于概念模型。E-R 图有三个基本成分：
\begin{itemize}[leftmargin=1.6em, itemsep=1pt]
  \item 实体（型）：用\textbf{矩形}表示，是客观存在并可以相互区分的事物；
  \item 属性：用\textbf{椭圆形}表示，与相应实体（或联系）用无向边相连；属性中加下划线的是码属性；
  \item 联系：用\textbf{菱形}表示，与有关实体用无向边相连，联系本身也可以有属性。
\end{itemize}

\textbf{（2）概念模型（E-R 模型）的作用（2 分）}

1）概念模型是现实世界到机器世界的一个中间层次，它能够真实、充分地反映现实世界中事物及其相互联系，且与具体的 DBMS 无关；

2）E-R 图直观易懂，便于数据库设计人员与用户之间交流意见，是需求分析成果的规范化表达；

3）概念模型独立于具体的 DBMS 和数据的物理存储结构，便于向关系、层次、网状等各种数据模型转换；

4）概念结构设计阶段先设计各局部应用的 E-R 图（分 E-R 图），再合并、优化为全局 E-R 图；在逻辑结构设计阶段，按转换规则（一个实体型转换为一个关系模式；1:1 与 1:n 联系可与任意一端对应的关系模式合并；m:n 联系必须单独转换为一个关系模式）把 E-R 图转换为关系模式。

\textbf{（3）联系的三种类型（1 分）}

两个实体型之间的联系可分为：一对一联系（1:1）、一对多联系（1:n）、多对多联系（m:n）。
\end{analysis}

\q 数据库恢复的基本原理、故障种类与恢复方法：
\begin{analysis}
\textbf{（1）基本原理（1.5 分）}

数据库恢复的基本原理是\textbf{利用冗余}。也就是说，数据库中任何一部分被破坏或丢失的数据，都可以利用存储在系统其他地方的冗余数据来重建。冗余数据主要通过两种途径建立：一是\textbf{数据转储（备份）}，即 DBA 定期把整个数据库复制到磁带或另一块磁盘上保存；二是\textbf{登记日志文件}，即把事务对数据库的每一次更新操作登记到日志文件中。恢复系统就是利用转储得到的后备副本和日志文件中的记录，把数据库恢复到某个已知的正确状态。

\textbf{（2）故障种类与相应的恢复方法（3.5 分）}

1）\textbf{事务故障}：某个事务在运行过程中未能正常终结（非预期、非正常的终止，如运算溢出、死锁被选为牺牲者、违反完整性约束等）。恢复方法：反向扫描日志文件，对该事务已经执行的更新操作执行 UNDO（撤销），把数据恢复到修改前的值。

2）\textbf{系统故障（软故障）}：造成系统停止运转、需要重新启动的任何事件，如 CPU 故障、操作系统崩溃、突然停电等。此时内存中的信息（包括数据库缓冲区中的内容）全部丢失，但磁盘上的数据库未受损。恢复方法：正向扫描日志文件，找出故障发生前已提交的事务放入 REDO 队列，找出故障发生时尚未提交的事务放入 UNDO 队列；对 UNDO 队列中的事务执行 UNDO，对 REDO 队列中的事务执行 REDO，从而把数据库恢复到一致状态。为减少 REDO 的工作量，可采用检查点（checkpoint）技术，在日志中周期性地记录检查点，只需从最近一个检查点开始恢复。

3）\textbf{介质故障（硬故障）}：外存被破坏，如磁盘损坏、磁头碰撞、瞬时强磁场干扰等，破坏性最大。恢复方法：首先重装数据库，即装入最新的数据库后备副本，使数据库恢复到最近一次转储时的状态；然后装入相应的日志文件副本，重做（REDO）自上次转储以来所有已提交的事务，把数据库恢复到故障发生前的一致状态。

（补充：为提高可靠性，重要的数据库系统通常采用双机热备份、磁盘镜像或 RAID 等硬件冗余措施。）
\end{analysis}

% ============================================================
\answerpart{四、综合题（2 题，每题 5 分，共 10 分）}
% ============================================================

\q SQL 题参考答案：

\textbf{（1）定义表 Work（1 分）}
\begin{lstlisting}[language=SQL]
CREATE TABLE Work (
    Eno   CHAR(6),
    Pno   CHAR(5),
    Hours INT,
    PRIMARY KEY (Eno, Pno),
    FOREIGN KEY (Eno) REFERENCES Emp(Eno) ON DELETE CASCADE,
    FOREIGN KEY (Pno) REFERENCES Project(Pno) ON DELETE CASCADE
);
\end{lstlisting}

\textbf{（2）查询“研发部”中工资高于本部门平均工资的职工（2 分）}
\begin{lstlisting}[language=SQL]
SELECT E.Eno, E.Ename, E.Salary
FROM Emp E, Dept D
WHERE E.Dno = D.Dno
  AND D.Dname = '研发部'
  AND E.Salary > (SELECT AVG(E2.Salary)
                  FROM Emp E2
                  WHERE E2.Dno = E.Dno)
ORDER BY E.Salary DESC;
\end{lstlisting}

\textbf{（3）查询参与 2 个及以上项目且总工时超过 100 小时的职工（2 分）}
\begin{lstlisting}[language=SQL]
SELECT E.Eno, E.Ename, COUNT(*) AS 参与项目数, SUM(W.Hours) AS 总工时
FROM Emp E, Work W
WHERE E.Eno = W.Eno
GROUP BY E.Eno, E.Ename
HAVING COUNT(*) >= 2 AND SUM(W.Hours) > 100
ORDER BY SUM(W.Hours) DESC;
\end{lstlisting}

\begin{analysis}
\textbf{第（1）问：}\code{PRIMARY KEY (Eno, Pno)} 定义复合主码（表级约束）；两个外码分别用 \code{FOREIGN KEY ... REFERENCES ...} 定义，题目要求的级联删除用 \code{ON DELETE CASCADE} 实现——当 \code{Emp} 或 \code{Project} 中的元组被删除时，\code{Work} 中对应的元组自动被删除，从而保证参照完整性。

\textbf{第（2）问：}部门名 \code{Dname} 在 \code{Dept} 中，而工资 \code{Salary}、部门号 \code{Dno} 在 \code{Emp} 中，因此需要 \code{Emp} 与 \code{Dept} 连接。内层子查询 \code{SELECT AVG(E2.Salary) FROM Emp E2 WHERE E2.Dno = E.Dno} 引用了外层查询的 \code{E.Dno}，属于\textbf{相关子查询}：对外层查询的每一个元组都要执行一次，求出该职工所在部门的平均工资，再与本人的工资比较。“高于”应写成严格大于 \code{>}（若写 \code{>=} 则把等于平均工资的人也选出来了）。结果按工资降序用 \code{ORDER BY E.Salary DESC}。

\textbf{第（3）问：}职工信息在 \code{Emp} 中、工时在 \code{Work} 中，需要两表连接后按职工分组（\code{GROUP BY E.Eno, E.Ename}）。“参与 2 个及以上项目”对应 \code{COUNT(*) >= 2}，“项目总工时超过 100 小时”对应 \code{SUM(W.Hours) > 100}，二者都是对分组结果的筛选，必须写在 \code{HAVING} 中；\code{SELECT} 列表中的非分组列只有出现在 \code{GROUP BY} 中的 \code{E.Eno}、\code{E.Ename} 是合法的。排序用 \code{ORDER BY SUM(W.Hours) DESC}。

\textbf{共性易错点：}组函数（\code{COUNT}、\code{SUM}、\code{AVG} 等）不能出现在 \code{WHERE} 子句中；\code{WHERE} 在分组前筛选元组，\code{HAVING} 在分组后筛选分组；\code{COUNT(*)} 统计所有行，\code{COUNT(列名)} 忽略该列为空值的行。
\end{analysis}

\q 规范化题参考答案：

\textbf{（1）候选码与主属性（1.5 分）}

$R$ 有\textbf{两个}候选码：\textbf{(Sno, Cno, Tno)} 和 \textbf{(Sno, Cno, Room)}。

推导（闭包运算）：
\begin{itemize}[leftmargin=1.6em, itemsep=1pt]
  \item 单属性闭包：Sno$^{+}=\{$Sno$\}$，Cno$^{+}=\{$Cno$\}$，Tno$^{+}=\{$Tno, Tname$\}$，Room$^{+}=\{$Room$\}$；
  \item (Sno, Cno)$^{+}=\{$Sno, Cno, Tno, Tname$\}$（由 (Sno,Cno)$\to$Tno，再由 Tno$\to$Tname），不含 Room，故 (Sno, Cno) 不是候选码；
  \item (Sno, Cno, Tno)$^{+}=\{$Sno, Cno, Tno, Tname$\}\cup\{$Room$\}=R$ 的全部属性（由 (Tno,Cno)$\to$Room，Tno 与 Cno 均在此集合中），故 (Sno, Cno, Tno) 是候选码；其真子集 (Sno,Cno)、(Sno,Tno)、(Cno,Tno) 的闭包分别为 $\{$Sno,Cno,Tno,Tname$\}$、$\{$Sno,Tno,Tname$\}$、$\{$Cno,Tno,Tname,Room$\}$，都推不出全部属性，故它是最小的；
  \item (Sno, Cno, Room)$^{+}=\{$Sno, Cno, Room, Tno, Tname$\}=R$ 的全部属性，故 (Sno, Cno, Room) 也是候选码；其真子集 (Sno,Cno)、(Sno,Room)、(Cno,Room) 均推不出全部属性，故它也是最小的。
\end{itemize}
（也可这样判断：$F$ 中任何函数依赖的右部都不含 Sno 和 Cno，故 Sno、Cno 必须包含在每个候选码中；而 (Sno,Cno) 推不出 Room，要得到 Room 只能补充 Tno（由 (Tno,Cno)$\to$Room）或直接补充 Room，因此候选码恰为上述两个。）

主属性：\textbf{Sno、Cno、Tno、Room}；非主属性：\textbf{Tname}。

\textbf{（2）$R$ 最高属于第几范式（1.5 分）}

$R$ 最高属于 \textbf{1NF}。

理由：$R$ 的所有属性都是不可再分的原子属性，故满足 1NF。候选码 (Sno, Cno, Tno) 的真子集 Tno 能决定非主属性 Tname，即存在
\[
\text{Tno}\to\text{Tname},
\]
其中 Tno 是候选码 (Sno, Cno, Tno) 的真子集，Tname 是非主属性，这是\textbf{非主属性对码的部分函数依赖}，违反 2NF，因此 $R$ 不属于 2NF，更不属于 3NF。故 $R$ 最高为 1NF。

\textbf{（3）分解为 3NF 并用 Chase 算法验证无损连接（2 分）}

\textbf{第一步，求最小依赖集。} $F$ 中各函数依赖的右部都是单属性；左部没有多余属性（如 (Sno,Cno)$\to$Tno 中去掉 Sno 或 Cno 后都推不出 Tno，(Tno,Cno)$\to$Room 中去掉 Tno 或 Cno 后也都推不出 Room）；逐个检查冗余：去掉 (Sno,Cno)$\to$Tno 后 (Sno,Cno)$^{+}=\{$Sno,Cno$\}$，不含 Tno，不冗余；去掉 Tno$\to$Tname 后 Tno$^{+}=\{$Tno$\}$，不冗余；去掉 (Tno,Cno)$\to$Room 后 (Tno,Cno)$^{+}=\{$Tno,Cno,Tname$\}$，不含 Room，不冗余。因此 $F_{\min}=F$。

\textbf{第二步，按保持函数依赖的 3NF 分解算法构造分解。} 对每个依赖 $X\to A$ 取关系模式 $XA$，得
\[
\rho=\{R_1,R_2,R_3\},
\]
\[
R_1(\text{Sno},\text{Cno},\text{Tno}),\quad
R_2(\text{Tno},\text{Tname}),\quad
R_3(\text{Tno},\text{Cno},\text{Room}).
\]
三个模式互不包含，不能去掉任何一个；而且 $R_1$ 恰好就是 $R$ 的候选码 (Sno, Cno, Tno)，已经包含了候选码，因此不需要再补充候选码关系模式。

每个依赖都落在某个模式中：(Sno,Cno)$\to$Tno 在 $R_1$，Tno$\to$Tname 在 $R_2$，(Tno,Cno)$\to$Room 在 $R_3$，故分解\textbf{保持函数依赖}。各模式的码分别为 (Sno,Cno)、Tno、(Tno,Cno)，每个函数依赖的决定因素都是该模式的码，所以 $R_1,R_2,R_3$ 均属于 BCNF，当然都属于 3NF。

\textbf{第三步，用表格法（Chase 算法）验证无损连接性。}

初始表格：行 $r_1,r_2,r_3$ 分别对应 $R_1,R_2,R_3$；若属性 $A_j$ 在 $R_i$ 中则填 $a_j$，否则填 $b_{ij}$。

\begin{center}
\begin{tabular}{lccccc}
\toprule
 & Sno & Cno & Tno & Tname & Room \\
\midrule
$r_1$（$R_1$） & $a_1$ & $a_2$ & $a_3$ & $b_{14}$ & $b_{15}$ \\
$r_2$（$R_2$） & $b_{21}$ & $b_{22}$ & $a_3$ & $a_4$ & $b_{25}$ \\
$r_3$（$R_3$） & $b_{31}$ & $a_2$ & $a_3$ & $b_{34}$ & $a_5$ \\
\bottomrule
\end{tabular}
\end{center}

依次使用 $F$ 中的函数依赖修改表格：
\begin{itemize}[leftmargin=1.6em, itemsep=1pt]
  \item 用 (Sno,Cno)$\to$Tno：没有任何两行的 Sno、Cno 同时相同（$r_1$ 的 Sno $=a_1$，$r_2$、$r_3$ 的 Sno 分别是 $b_{21}$、$b_{31}$），故表格不发生变化；
  \item 用 Tno$\to$Tname：$r_1,r_2,r_3$ 的 Tno 列都是 $a_3$，属同一组，须把 Tname 列取值统一；其中 $r_2$ 的 Tname 为区分符号 $a_4$，故把 $r_1$ 的 $b_{14}$、$r_3$ 的 $b_{34}$ 都改为 $a_4$；
  \item 用 (Tno,Cno)$\to$Room：Tno 列与 Cno 列都相同的行是 $r_1$（Cno $=a_2$）与 $r_3$（Cno $=a_2$），须把它们的 Room 列统一；$r_3$ 的 Room 为区分符号 $a_5$，故把 $r_1$ 的 $b_{15}$ 改为 $a_5$；
  \item 此时 $r_1$ 已变为 $(a_1,a_2,a_3,a_4,a_5)$，该行全为区分符号，说明表格中存在一行全为 $a$ 的行。
\end{itemize}

按 Chase 算法，若修改后的表格中存在一行全为区分符号，则分解 $\rho$ 具有\textbf{无损连接性}。因此 $\rho=\{R_1,R_2,R_3\}$ 既保持函数依赖又具有无损连接性，且每个模式都属于 3NF（实际上都属于 BCNF），满足题目要求。

\begin{analysis}
\textbf{易错点提示：}
1）本题的难点是候选码不唯一。(Sno, Cno) 不是候选码，因为 (Sno,Cno)$^{+}$ 中推不出 Room；只有再补上 Tno 或 Room 才能得到全部属性，由此得到两个候选码。
2）判断 2NF 要用“非主属性是否部分依赖于\textbf{某个}候选码”，只要存在一处即违反。Tno$\to$Tname 中 Tno 是候选码 (Sno,Cno,Tno) 的真子集，Tname 又是非主属性，所以只达到 1NF。
3）无损连接的判定必须给出依据。本答案采用 Chase 表格法；也可以用二元分解判定定理逐步连接：$R_1\cap R_2=\{$Tno$\}$，$R_2-R_1=\{$Tname$\}$，而 Tno$\to$Tname $\in F^{+}$，故 $R_1$ 与 $R_2$ 的连接无损，连接结果 $T_{12}(\text{Sno},\text{Cno},\text{Tno},\text{Tname})$；再令 $T_{12}\cap R_3=\{$Tno, Cno$\}$，$R_3-T_{12}=\{$Room$\}$，而 (Tno,Cno)$\to$Room $\in F^{+}$，故与 $R_3$ 的连接也无损。两种方法结论一致。
4）不能把 $R_3$ 合并到 $R_1$ 中得到 $(\text{Sno},\text{Cno},\text{Tno},\text{Room})$：该模式的码为 (Sno, Cno)，而其中的函数依赖 (Tno, Cno)$\to$Room 的左部 (Tno, Cno) 不是码，右部 Room 又是非主属性，不满足 3NF 的条件，所以合并后的模式达不到 3NF；同时去掉 $R_3$ 还会使 (Tno,Cno)$\to$Room 丢失，不再保持函数依赖。
\end{analysis}

\end{document}
